\section{Protocols Under Comparison}
\label{sec:protocols}

\begin{table*}[!t]
\caption{The ten configurations compared. ``Control model'' determines which setup messages are charged: distributed protocols pay advertisement, join and schedule; the chain protocol pays a hop-by-hop token; centralized protocols pay only the downlink schedule, their uplink state being piggybacked on data traffic (Section~\ref{sec:validity}).}
\label{tab:taxonomy}
\centering
\begin{tabular}{llllll}
\toprule
Protocol & Gen. & Head-selection rule & Control model & Re-cluster & Offline train \\
\midrule
Direct (no clustering) & -- & None & None & -- & No \\
LEACH & G1 & Randomised rotation & Distributed & Every round & No \\
PEGASIS & G1 & Greedy chain & Chain & On death & No \\
TEEN & G1 & LEACH + hard/soft threshold & Distributed & Every round & No \\
APTEEN & G1 & TEEN + count-down timer & Distributed & Every round & No \\
NSGA-II & G2 & Multi-objective search (3 obj.) & Centralized & Every 5 rounds & No \\
Fuzzy T2 & G2 & Interval type-2 rule base (27 rules) & Centralized & Every round & No \\
SOM & G3 & Self-organising map & Centralized & Every round & Yes \\
DQN & G3 & Learned $Q$-value ranking & Centralized & Every round & Yes (frozen) \\
GCN & G3 & Learned graph scoring & Centralized & Every round & Yes (frozen) \\
\bottomrule
\end{tabular}
\end{table*}


Nine protocols are implemented, spanning three generations, plus a
direct-transmission baseline in which every node sends its own reading straight
to the sink. Table~\ref{tab:taxonomy} sets out the ten configurations against
class, control-traffic model, re-clustering cadence and pre-training.

All nine share one structural contract. Each targets the same nominal cluster
count $k = \max(1, \mathrm{round}(p \cdot N_{\mathrm{alive}}))$, recomputed each
round so $k$ decays as the network dies, with membership always nearest-head by
Euclidean distance. What each protocol returns is a set of head identities and a
membership map, nothing else; energy, packet sizing, fusion charges, the channel
and reading accounting belong to the engine (Section~\ref{sec:invariants}).
Realized head counts are not forced to match the target, because clamping
LEACH's stochastic election would be implementing a different algorithm.

\subsection{Generation One: Heuristics}

LEACH~\cite{heinzelman2000leach} implements the original election threshold
unmodified: a node that has not yet been a head in the current epoch self-elects
with probability $T(n) = p/(1 - p(r \bmod 1/p))$, and the served set clears
every $1/p = 20$ rounds. The epoch mechanism turns out to be LEACH's real
strength, because it guarantees each surviving node serves exactly once per
epoch, close to optimal for first death specifically, and no other protocol here
has an equivalent guarantee.

PEGASIS~\cite{lindsey2002pegasis} abandons clusters for a greedy chain built
from the node farthest from the sink, rebuilt only when a death invalidates it,
with the leader rotating round-robin over chain positions. Data flows inward
from both ends and fuses at every hop, so each non-leader emits one packet and
the whole network's readings reach the sink in a single transmission. PEGASIS is
the sole multi-hop protocol here and the sole beneficiary of ideal fusion at
full scale.

TEEN~\cite{manjeshwar2001teen} subclasses LEACH and reuses its clustering
unchanged. This is the correct reading of the protocol rather than a shortcut:
TEEN's contribution is entirely the reporting rule, and holding clustering fixed
makes the comparison controlled. A node transmits only when its reading exceeds
a hard threshold and has moved by a soft threshold since its last transmission.
APTEEN~\cite{manjeshwar2002apteen} adds the periodic half, so a node that never
crosses the hard threshold still reports every 50 rounds. Only the
reporting-frequency mechanism is modeled; APTEEN's query facilities are out of
scope.

\subsection{Generation Two: Explicit Optimization}

NSGA-II~\cite{deb2002nsga2} poses head selection as a genuine multi-objective
problem, solved with the full algorithm, fast non-dominated sort, crowding
distance, binary tournament selection and elitist $(\mu + \lambda)$ survival,
rather than a weighted sum standing in for it. A genome is a vector of $k$
distinct living node indices, so cardinality is fixed by construction. Three
objectives are minimized: total energy of one round under the candidate
assignment, the negated minimum residual energy among selected heads, and the
standard deviation of cluster membership counts. The third was chosen against
the more common intra-cluster distance variance, because distance already
dominates the first objective through the $d^2$ and $d^4$ terms. The first
objective calls the engine's own radio functions, so the search optimizes
exactly the cost model it will later be charged against, and selection from the
final front is by knee point, which introduces no hand-tuned weights that could
act as a back door to calibrating against published results. Evolutionary and
swarm formulations have a long history~\cite{ferentinos2007ga,kulkarni2011pso};
what is new is not the formulation but that it is charged against the same
engine as everything else.

Interval type-2 fuzzy selection scores every living node from residual energy,
distance to sink and local density, with Gaussian membership functions of
uncertain mean, a 10\% footprint of uncertainty, Karnik-Mendel type reduction
and midpoint
defuzzification~\cite{karnik2001centroid,wu2009ekm,mendel2002type2}. The 27-rule
base is generated by $\mathrm{score} = 2E + (2 - D) + D_n$ over three-level
inputs, giving a table monotone in every input by construction, which is what
makes it defensible without tuning; the double weight on residual energy is a
designer's choice and we label it as one. Fuzzy head election has an established
line of work~\cite{gupta2005fuzzy,kim2008chef,bagci2013eaucf}, and our
contribution is not the rule base but its measurement under the same harness as
everything else.

\subsection{Generation Three: Learned Selection}

SOM~\cite{kohonen1990som} trains a Kohonen map with $k$ units over normalized
position and residual energy, retrained from scratch on the same five-round
cadence as NSGA-II and the fuzzy system, and the node nearest each converged
unit becomes a head.

DQN~\cite{mnih2015dqn,watkins1992qlearning} and GCN~\cite{kipf2017gcn} are
deliberately built to isolate a single difference. Both consume the identical
four normalized features, residual energy, distance to sink, local density and
rounds since last serving as head, with the GCN importing the DQN's feature
function rather than redefining it. Both use hand-written forward and backward
passes verified against central finite differences, and both are trained on
seeds 100 to 129, frozen, and evaluated on the disjoint seeds 0 to 29. Only the
architecture and the training signal differ.

The DQN is a $4 \rightarrow 16 \rightarrow 2$ multilayer perceptron trained with
experience replay, a synchronized target network, Adam~\cite{kingma2015adam} and
$\gamma = 0.999$. At evaluation, rather than letting each node act on its own
greedy action, which would make the head count fluctuate between zero and the
entire population, nodes are ranked by the advantage
$Q_{\mathrm{CH}} - Q_{\overline{\mathrm{CH}}}$ and the top $k$ selected. This
makes it a learned ranking function rather than a set of independent agents, a
real narrowing of the reinforcement-learning framing~\cite{luong2019drl} that we
document rather than gloss over.

The GCN is two graph-convolutional layers of width 16 over a symmetrized
6-nearest-neighbor graph, trained self-supervised against a differentiable proxy
for energy balance. It is deliberately not trained to imitate a heuristic, since
that would make any comparison against that heuristic vacuous, and shares no
code path with the DQN's training loop. Selection is greedy top-$k$ subject to a
15~m minimum separation.

The contrast this buys is the paper's cleanest: with identical inputs, the DQN
has temporal credit assignment and the GCN does not. The GCN can see rotation
state through its fourth feature, but nothing in its single-round loss rewards
acting on it. Section~\ref{sec:credit} measures what that costs.

\subsection{Disclosure on the DQN's Hyperparameters}
\label{sec:dqn-disclosure}

One protocol's parameters were revised after observing a failure, and we state
it rather than let it be inferred~\cite{henderson2018matters}. The initial
$\gamma = 0.95$ gave an effective horizon of about twenty rounds against a
first-death timescale three orders of magnitude longer, and it diverged. The fix
raised $\gamma$ to 0.999, normalized the return magnitude and extended training
from 30 to 100 episodes, leaving the reward function unchanged. A stopping rule
was then fixed in advance, recording the result as final regardless of outcome,
after which it beat random rotation on all five held-out seeds. Stating the
pre-commitment is what makes that number credible, since without it ``we fixed
the DQN until it worked'' and ``we fixed a divergence and stopped'' are
indistinguishable from the outside.
